\section{Introduction}
Task-specific vision models offer efficient predictions, while general vision--language models (VLMs) can handle a broader range of inputs at greater computational cost. Calling a VLM only when needed may improve accuracy without paying this cost on every input, a motivation also reflected in prompt-routing services~\cite{bedrockrouting}. We consider a system in which a specialist first predicts a class, a set of boxes or an event label. A router then decides whether to keep that answer or replace it with the VLM's answer.

A natural rule is to call the VLM when the specialist has low confidence. However, uncertainty alone does not show whether the replacement will help. The VLM may fail on the same example, or overturn an answer that was already correct. Even a VLM with lower average accuracy can be useful if its strengths complement those of the specialist. Routing must therefore learn where the two models differ, rather than only where the specialist is uncertain.

Prior work offers several starting points. Confidence estimation~\cite{calibration} and selective prediction~\cite{selectivenet} identify uncertain inputs. Learning to defer also accounts for an imperfect expert~\cite{defer,calibrateddefer}. Its joint-classification formulations learn class outputs alongside deferral, changing the retained predictor and requiring adaptation for box or event outputs. Post-hoc deferral~\cite{posthocdefer} keeps an existing predictor fixed, making it a closer comparison for answer replacement. We adapt its loss to our router inputs and backbone. Text routing and cascading~\cite{frugal,routellm,cascaderouting} and multimodal routing~\cite{vlrouterbench,mmrbench,relope,streamsense} address related choices, but their model interfaces and evaluation settings require adaptation here. Speculative decoding~\cite{specdecode,specsampling} verifies draft tokens rather than choosing between completed visual answers.

We introduce \textbf{Learned Routing (LR)}, a compact router that uses input features, specialist representations, predictions and confidence. During training, answers from both models reveal when replacement helps or hurts. At inference, LR uses only information available before the VLM call. We distinguish what the router learns from how its outputs rank inputs: the same trained model can produce different rankings when its score is changed. Our default normalizes estimated replacement gains, choosing the normalization strength on validation data.

Our contributions are as follows.

\noindent\textbf{(1) Unified evaluation.} We bring confidence routing, LR and post-hoc deferral~\cite{posthocdefer} into cross-task experiments spanning fine-grained recognition, grounding, two-image reasoning and safety-event recognition. Across 16 model pairs, full-call-range curves show where collaboration improves performance and where the specialist alone is preferable.

\noindent\textbf{(2) Routing method and empirical finding.} We develop a compact router with a normalized ranking score, and compare training objectives separately from score construction. In some tasks, changing the score without retraining has a larger effect than changing among the tested objectives. The preferred score remains task-dependent.

\noindent\textbf{(3) Open-source release.} We release all experiment code and per-example experimental data, the 80 router checkpoints used in the main table, and the configurations, prompts and reproduction scripts. These materials support retraining the routers, reconstructing the reported curves and comparing additional methods.
